% xystyle.tex — XY 스테이지 다이어그램 공통 스타일 (% xystyle.tex — XY 스테이지 다이어그램 공통 스타일 (% xystyle.tex — XY 스테이지 다이어그램 공통 스타일 (% xystyle.tex — XY 스테이지 다이어그램 공통 스타일 (\input{xystyle})
%   흰 바탕, 직각 상자, 회색 + 남색 1색 강조, 수식은 LaTeX (Computer Modern), 한글은 kotex
\usepackage{kotex}
\usepackage{amsmath,amssymb}
\usetikzlibrary{arrows.meta,positioning,calc,fit,backgrounds}
\definecolor{navy}{HTML}{1F3A5F}
\definecolor{navybg}{HTML}{E9EEF5}
\definecolor{ink}{HTML}{333333}
\definecolor{hwbg}{HTML}{F3F3F3}
\definecolor{mute}{HTML}{777777}
\tikzset{
  >={Stealth[length=2.0mm,width=1.6mm]},
  every node/.style={text=ink},
  blk/.style={draw=ink, line width=0.6pt, fill=white, align=center, inner sep=4pt,
              minimum height=1.0cm, font=\small},
  key/.style={blk, draw=navy, fill=navybg, line width=0.8pt},
  hw/.style={blk, fill=hwbg},
  opt/.style={blk, dashed, draw=mute},
  sig/.style={->, line width=0.6pt, draw=ink},
  optsig/.style={->, line width=0.6pt, dashed, draw=mute},
  lab/.style={font=\footnotesize, align=center, inner sep=1.5pt},
  note/.style={font=\footnotesize, align=left, text=mute},
  grp/.style={draw=mute, dashed, line width=0.5pt, inner sep=7pt},
  ttl/.style={font=\footnotesize\bfseries, text=navy, anchor=south west, inner sep=1pt},
}
\newcommand{\tp}[1]{\texttt{\footnotesize #1}}
)
%   흰 바탕, 직각 상자, 회색 + 남색 1색 강조, 수식은 LaTeX (Computer Modern), 한글은 kotex
\usepackage{kotex}
\usepackage{amsmath,amssymb}
\usetikzlibrary{arrows.meta,positioning,calc,fit,backgrounds}
\definecolor{navy}{HTML}{1F3A5F}
\definecolor{navybg}{HTML}{E9EEF5}
\definecolor{ink}{HTML}{333333}
\definecolor{hwbg}{HTML}{F3F3F3}
\definecolor{mute}{HTML}{777777}
\tikzset{
  >={Stealth[length=2.0mm,width=1.6mm]},
  every node/.style={text=ink},
  blk/.style={draw=ink, line width=0.6pt, fill=white, align=center, inner sep=4pt,
              minimum height=1.0cm, font=\small},
  key/.style={blk, draw=navy, fill=navybg, line width=0.8pt},
  hw/.style={blk, fill=hwbg},
  opt/.style={blk, dashed, draw=mute},
  sig/.style={->, line width=0.6pt, draw=ink},
  optsig/.style={->, line width=0.6pt, dashed, draw=mute},
  lab/.style={font=\footnotesize, align=center, inner sep=1.5pt},
  note/.style={font=\footnotesize, align=left, text=mute},
  grp/.style={draw=mute, dashed, line width=0.5pt, inner sep=7pt},
  ttl/.style={font=\footnotesize\bfseries, text=navy, anchor=south west, inner sep=1pt},
}
\newcommand{\tp}[1]{\texttt{\footnotesize #1}}
)
%   흰 바탕, 직각 상자, 회색 + 남색 1색 강조, 수식은 LaTeX (Computer Modern), 한글은 kotex
\usepackage{kotex}
\usepackage{amsmath,amssymb}
\usetikzlibrary{arrows.meta,positioning,calc,fit,backgrounds}
\definecolor{navy}{HTML}{1F3A5F}
\definecolor{navybg}{HTML}{E9EEF5}
\definecolor{ink}{HTML}{333333}
\definecolor{hwbg}{HTML}{F3F3F3}
\definecolor{mute}{HTML}{777777}
\tikzset{
  >={Stealth[length=2.0mm,width=1.6mm]},
  every node/.style={text=ink},
  blk/.style={draw=ink, line width=0.6pt, fill=white, align=center, inner sep=4pt,
              minimum height=1.0cm, font=\small},
  key/.style={blk, draw=navy, fill=navybg, line width=0.8pt},
  hw/.style={blk, fill=hwbg},
  opt/.style={blk, dashed, draw=mute},
  sig/.style={->, line width=0.6pt, draw=ink},
  optsig/.style={->, line width=0.6pt, dashed, draw=mute},
  lab/.style={font=\footnotesize, align=center, inner sep=1.5pt},
  note/.style={font=\footnotesize, align=left, text=mute},
  grp/.style={draw=mute, dashed, line width=0.5pt, inner sep=7pt},
  ttl/.style={font=\footnotesize\bfseries, text=navy, anchor=south west, inner sep=1pt},
}
\newcommand{\tp}[1]{\texttt{\footnotesize #1}}
)
%   흰 바탕, 직각 상자, 회색 + 남색 1색 강조, 수식은 LaTeX (Computer Modern), 한글은 kotex
\usepackage{kotex}
\usepackage{amsmath,amssymb}
\usetikzlibrary{arrows.meta,positioning,calc,fit,backgrounds}
\definecolor{navy}{HTML}{1F3A5F}
\definecolor{navybg}{HTML}{E9EEF5}
\definecolor{ink}{HTML}{333333}
\definecolor{hwbg}{HTML}{F3F3F3}
\definecolor{mute}{HTML}{777777}
\tikzset{
  >={Stealth[length=2.0mm,width=1.6mm]},
  every node/.style={text=ink},
  blk/.style={draw=ink, line width=0.6pt, fill=white, align=center, inner sep=4pt,
              minimum height=1.0cm, font=\small},
  key/.style={blk, draw=navy, fill=navybg, line width=0.8pt},
  hw/.style={blk, fill=hwbg},
  opt/.style={blk, dashed, draw=mute},
  sig/.style={->, line width=0.6pt, draw=ink},
  optsig/.style={->, line width=0.6pt, dashed, draw=mute},
  lab/.style={font=\footnotesize, align=center, inner sep=1.5pt},
  note/.style={font=\footnotesize, align=left, text=mute},
  grp/.style={draw=mute, dashed, line width=0.5pt, inner sep=7pt},
  ttl/.style={font=\footnotesize\bfseries, text=navy, anchor=south west, inner sep=1pt},
}
\newcommand{\tp}[1]{\texttt{\footnotesize #1}}
